\section {Printing a linked list in reverse (Algorithms)}

Print a singly linked list of $N$ elements from its end to its start, using only $O(1)$ memory.
You may not change the list. Make the running time as small as you can.

\subsection{Solution}

The simplest solution is to walk from the head to the last element and print it, then walk from the head to the element before last and print it,
and so on. This takes $O(N^2)$ time. Let's call it the naive method; it prints any segment of $m$ elements backwards in $O(m^2)$ time.

\textbf{First improvement.} Split the list into $\sqrt{N}$ blocks of $\sqrt{N}$ elements each.
Go over the blocks from the last to the first. For each block, walk from the head to the start of the block, which takes $O(N)$,
and print the block backwards with the naive method, which takes $O((\sqrt{N})^2) = O(N)$.
There are $\sqrt{N}$ blocks, so the total time is $O(\sqrt{N} \cdot N) = O(N^{1.5})$.

\textbf{Second improvement.} Split the list into $N^{1/3}$ blocks of $N^{2/3}$ elements each,
and print each block backwards with the first improvement instead of the naive method.
A block of $m = N^{2/3}$ elements takes $O(m^{1.5}) = O(N)$ time to print, and walking to the start of the block also takes $O(N)$.
There are $N^{1/3}$ blocks, so the total time is $O(N^{4/3})$.

\textbf{And so on.} Repeating the trick $k$ times gives a running time of $O(N^{1+1/k})$.
Every level of the trick needs only a few variables, so for any fixed $k$ the memory is still $O(1)$,
and the running time is $O(N^{1+\epsilon})$ for any $\epsilon > 0$ we like.

If we allow $O(\log N)$ memory, we can go all the way: split the list into two halves, print the second half backwards and then the first half backwards,
recursively. Every level of the recursion walks over the list once, and there are $\log N$ levels, so this takes $O(N \log N)$ time.
